
%https://docs.arduino.cc/tutorials/nicla-vision/nicla-vision-imu



\section{Verwendung der Inertialmesseinheit (IMU) des Nicla Visions}

Eine \ac{imu} ist ein Bauteil, das aus verschiedenen Sensoren besteht, die Daten wie spezifische Kraft, Winkelgeschwindigkeit und Orientierung aufzeichnen. Die \ac{imu}  des Nicla Visions ist ein Baustein vom Type LSM6DSOXTR. Dieser besitzt ein Gyroskop und einen Beschleunigungsmesser.  Auf dem Bild \ref{IMUpositon} ist markiert, wo sich die \ac{imu} auf der Platine befindet.

Im Folgenden wird gezeigt, wie man innerhalb der Arduino IDE auf das Gyroskop und den Beschleunigungsmesser zugreift. Dazu wird die  Bibliothek \PYTHON{Arduino\_LSMDS63}. Die Werte können dann über den seriellen Monitor angeschaut werden.

Dazu wird eine Anwendung vorgestellt, die die Kräfte darstellt, die auf das Board wirken. DEr Beschleunigungssensor wird als \glqq Libelle\grqq{} verwenden, die Informationen über die Position des Brettes liefert. Mit dieser Anwendung ist es möglich, die relative Position des Boards abzulesen, ebenso wie die Gradzahl, wenn das Board nach oben, unten, links oder rechts gekippt wird. Die Ergebnisse werden über den seriellen Monitor sichtbar sein.






\begin{figure}
	\centering
	
	\includegraphics[width=0.5\textwidth]{IMU/nicla-vision-imu.png}
	
	\caption{Position der \ac{imu} auf dem Arduino Nicla Vision Board}\label{IMUpositon}
\end{figure}


%todo Ausarbeiten
%\subsection{Inertialmesseinheit (IMU)}


\subsection{Beschleunigungsmesser und Gyroskop}

Ein Beschleunigungsmesser ist ein elektromechanisches Gerät, das zur Messung von Beschleunigungskräften verwendet wird. Solche Kräfte können statisch sein, wie die kontinuierliche Schwerkraft, oder, wie bei vielen mobilen Geräten, dynamisch, um Bewegungen oder Vibrationen zu erfassen.


\begin{figure}
    \centering
    
    \includegraphics[width=0.8\textwidth]{IMU/nicla_vision_acceleration}
    
    \caption{Achsen des Beschleunigungsmessers des Nicla Visions.}\label{IMUaccelerometer}
\end{figure}

Ein Gyroskopsensor ist ein Gerät, das die Ausrichtung und Winkelgeschwindigkeit eines Objekts messen und beibehalten kann. Gyroskope sind fortschrittlicher als Beschleunigungsmesser, da sie die Neigung und seitliche Ausrichtung eines Objekts messen können, während ein Beschleunigungsmesser nur die lineare Bewegung messen kann. Gyroskopsensoren werden auch als ``Winkelgeschwindigkeitssensoren'' bezeichnet. Gemessen in Grad pro Sekunde, ist die Winkelgeschwindigkeit die Änderung des Drehwinkels des Objekts pro Zeiteinheit.


\begin{figure}
    \centering
    
    \includegraphics[width=0.8\textwidth]{IMU/nicla_vision_gyroscope.png}
    
    \caption{Achsen des Gyroskops des Nicla Visions.}\label{IMUgyroscope}
\end{figure}



\subsection{Bibliothek ArduinoLSM6DSOX}


Die Bibliothek ArduinoLSM6DSOX ermöglicht es, die auf dem Arduino-Board verfügbare \ac{imu}  zu verwenden. Die \ac{imu} ist ein LSM6DSOX, ein 3-Achsen-Beschleunigungsmesser und ein 3-Achsen-Gyroskop. Die \ac{imu} ist über I\textsuperscript{2}C\index{I\textsuperscript{2}C} mit dem Mikrocontroller verbunden. Die zurückgegebenen Werte sind vorzeichenbehaftete Fließkommazahlen.

\bigskip

Zur Verwendung der Bibliothek muss sie wie folgt in den Sketch  eingebunden werden:

\bigskip

\PYTHON{\#include <Arduino\_LSM6DSOX.h>}

\bigskip

Durch die Verwendung der Bibliothek wird der Sensor initialisiert.  Dabei werden folgende Auflösungen und Übertragungsrate eingehalten:


\begin{itemize}
  \item Der Bereich des Beschleunigungsmessers ist auf $\pm 4 g$ mit einer Auflösung von 0,122 mg eingestellt.
  \item Die Reichweite des Gyroskops ist auf $\pm 2000 dps$ mit einer Auflösung von 70 mdps eingestellt.
  \item Beschleunigungsmessers und Gyroskop haben eine feste Ausgangsdatenrate von 104 Hz.
\end{itemize}


%\subsubsection{Einrichten der Arduino IDE}
%
%Stellen Sie sicher, dass der neueste Nicla Core in der Arduino IDE installiert ist. Werkzeuge > Platine > Platinenmanager.... Hier müssen wir nach den Arduino Mbed OS Nicla Boards suchen und sie installieren. Nun müssen wir die für die IMU benötigte Bibliothek installieren. Gehen Sie zu Tools > Manage libraries.., und suchen Sie nach \PYTHON{Arduino\_LSM6DS3} und installieren Sie es.

\subsection{Sketch zur Verwendung der IMU}

\subsubsection{Aufbau des Sketchs zur Verwendung der IMU}

Den vollständigen Sketch ist in \ref{IMUSketch} dargestellt. Dieser muss auf das Board geladen werden.

Zunächst wird die Bibliothek eingebunden. Danach wird für jeden Wert der \ac{imu} einen Variable angelegt. 

In der Standardroutine \PYTHON{setup} wird die Methode \PYTHON{IMU.begin()} aufgerufen, so dass die Bibliothek initialisiert wird. Wenn die \ac{imu} inistialisiert ist, können die Abtastraten der Sensoren überprüft werden. Durch Aufruf von \PYTHON{IMU.accelerationSampleRate()} und \PYTHON{IMU.gyroscopeSampleRate()} wird die Abtastrate des jeweiligen Sensors in Hz ausgelesen.


In der Standardroutine \PYTHON{loop} des Sketches werden mit \PYTHON{IMU.accelerationAvailable()} und \PYTHON{IMU.gyroscopeAvailable()} die Sensoren überprüft, um zu sehen, ob Daten für die \ac{imu}-Sensoren verfügbar sind. Anschließend wird die Methode \PYTHON{IMU.readAcceleration(Ax, Ay, Az)} aufgerufen, um den Beschleunigungsmesser zu lesen. Er gibt den Wert der x-, y- und z-Achse zurück und aktualisiert die Variablen \PYTHON{Ax}, \PYTHON{Ay} und \PYTHON{Az}. Analog geschieht dies nun für das Gyroskop. Die Werte werden dann formatiert und  im seriellen Monitor dargestellt. Die Daten werden mit einem Intervall von 500 Millisekunden gesendet. Dies kann durch Ändern der Zeile \PYTHON{delay(500)} am unteren Ende der Skizze angepasst werden.


Nach dem erfolgreichen Hochladen des Codes auf das Board muss der serielle Monitor geöffnet werden.  Sobald er geöffnet ist, werden die Daten angezeigt.


\begin{code}
    {\scriptsize
        \lstinputlisting[language=C++]{../Code/Arduino/IMU/IMU.ino}
    }
    \caption{Sketch zur Verwendung der IMU}\label{IMUSketch}
\end{code}




